\pagetopsection{ Traits}


\label{sec:traits}

\begin{multicols*}{2}

\subsection*{Martial traits}

% TODO: some flavour text about trading spellcasting for combat prowess.

\paragraph{Sacrifical protection} is a rare technique that is taught for those knights who are ready to ensure that their patron is safe no matter the cost.

During enemy's turn, if creature within the jump distance is targeted by an attack, you can shield them with your own body.

After \texttt{Focus}, \texttt{Stress} and \texttt{Deny} are applied -- you leap: choose exactly half dice, rounded up, from the attack dice. 
Treat them as if it is an attack against you. Remove those dice from the attack that you interfered with.

Pass a \texttt{Vigor Save} or become \texttt{Fatigued}.

\paragraph{Weapon mastery} manifests when a knight develops an uncanny fascination with their weapons. 
They treat them not as mere tools of the craft, but as partners or trustworthy friends.

That allows them to see beyond the material, beyond the function, but towards possibilities.

Each time you roll a maximum value on a weapon die, put a mark on the weapon it belongs to. 
When amount of marks matches its largest weapon die -- you obtain a level of mastery with this weapon. 

When you use \texttt{Focus} or \texttt{Stress} and choose to duplicate the dice that belong to your weapon -- for each mastery level
you can keep the dice that otherwise would be discarded due to accumulated \texttt{Advantages}.

Weapon mastery levels do not transfer between any weapons even if they are of the same type, e.g. replacing one shortsword with another does not allow to keep the mastery level.

\paragraph{Fluency of Combat} occasionally comes to those who never engaged with the world of magic.

During fights, you can spend your \texttt{Guard} to move in a strange dance of blade and fury. 
When deciding what dice to spend on Gambits, you can use dice with value below \texttt{4} as long as you compensate the difference by reducing your \texttt{Guard}.

\newcolumn

\paragraph{Deadly precision} of your strike is often more valuable then the fury and might put into it. 

Instead of attacking normally, you can pick one weapon and pass a \texttt{Vigor save}. 
If successful -- treat the attack as if you attack only with this weapon and rolled maximum value of all its weapon dice.
On fail treat as if they rolled zeros.

% \paragraph{Riposte} is a risky and potentially deadly \texttt{Feat}. If you choose to \emph{Riposte}, you can use it instead of \texttt{Deny Feat}. 
% Roll your weapon dice as if you are performing an attack. 
% Spend any amount of \texttt{Hope} or \texttt{Focus} you allocated for \emph{Defense}, keep one weapon dice per point spent.
% % Choose the number of them equal to \texttt{Focus} you allocated to \emph{Defense}, spending one point per dice. You can also spend Hope to add those dice.

% You can match any of your dice with lower dice in the attack roll against you. 
% Deal damage equal to the sum of pairwise differences between matched dice. Discard enemy's dice that were matched. 
% The damage of the enemy's attack after \emph{Riposte} ignores your \texttt{Guard}.

% Pass a \texttt{Spirit} save or become \texttt{Fatigued}.

\subsection*{Arcane traits}

\paragraph{Rethreading dominion} permits a whole new level of control over spellweaving threads.

When choosing dice to store during \emph{rethreading the arcane} you can combine those dice instead of storing them individually according to steps below.

\colbox{
\begin{enumerate}[nosep, leftmargin=*]
    \item Identify dice to combine among the dice chosen for rethreading
    \item Mark one of them as accumulation die
    \item Add all values among identified dice and set the accumulation die to that value
\end{enumerate}
If the accumulation die can't hold that value -- \emph{rethreading of the arcane} fails.
}

For example, if you chose to store two \texttt{d6}'s with values \texttt{2} and \texttt{3} you can instead store them as a single \texttt{d6} with a value \texttt{5} on it. 

On the other hand if you choose to store two \texttt{d4}'s with the same values rethreading fails as \texttt{d4} can't store value of \texttt{5}. 
If one of them is \texttt{d6} instead, it becomes possible again.

\end{multicols*}
